% Generated draft: requires mathematical and source review.
\documentclass[11pt]{article}
\usepackage[margin=25mm]{geometry}
\usepackage{fontspec}
\setmainfont{Latin Modern Roman}
\usepackage{amsmath,amssymb,mathtools,unicode-math}
\setmathfont{Latin Modern Math}
\usepackage{enumitem,needspace,xurl,hyperref}
\setlength{\parindent}{0pt}
\setlength{\parskip}{5pt}
\setlength{\emergencystretch}{3em}
\newcommand{\N}{\mathbb{N}}
\newcommand{\Z}{\mathbb{Z}}
\newcommand{\Q}{\mathbb{Q}}
\newcommand{\R}{\mathbb{R}}
\newcommand{\C}{\mathbb{C}}
\newcommand{\F}{\mathbb{F}}
\newcommand{\A}{\mathbb{A}}
\newcommand{\PP}{\mathbb{P}}
\newcommand{\E}{\mathbb{E}}
\newcommand{\eps}{\varepsilon}
\newcommand{\dd}{\,\mathrm{d}}
\begin{document}
\section*{UnsolvedMath ID: 20000543}
\textbf{Source problem code:} catalogue entry
\subsection{Definitions and mathematical statement}
Let \(M\) be a von Neumann factor, so its center is \(\C1\), and let \(A\subseteq M\) be a unital maximal abelian von Neumann subalgebra (MASA), equivalently \(A'\cap M=A\). A state means a positive unital functional on the underlying \(C^*\)-algebra, not necessarily a normal state. A pure state is an extreme point of the state space; for abelian \(A\), these are its characters \(\chi\in\widehat A\). Put
\[
 \operatorname{Ext}_M(\chi)=\{\varphi\in M^*:\varphi\text{ a state},\ \varphi|_A=\chi\}.
\]
The source asks which inclusions have the all-pure-state unique extension property
\[
 \#\operatorname{Ext}_M(\chi)=1\quad\text{for every }\chi\in\widehat A.
\]
Failure means at least one pure state has distinct extensions, not that every pure state fails uniqueness.

The November 2014 compilation reports: MSS gives uniqueness for discrete MASAs of \(B(H)\) on separable \(H\); any continuous part of such a MASA prevents the all-state property, as in the original Kadison--Singer paper; every MASA of a type II factor acting on separable \(H\) fails it; nonseparable examples can behave either way; and no type III results were known there. These are historical assessments, not present-day claims that no type III theorem exists.

The supplied later assessment states that for every MASA \(A\subseteq M\) with \(M\) having separable predual, the all-pure-state property holds exactly when \(M\) is type I and \(A\) is the range of a normal conditional expectation. It says, in particular, that every MASA in a separable-predual type II or III factor has a singular pure state with nonunique, necessarily singular, extensions, while the arbitrary-predual factor problem remains unclassified. These claims require their absent primary theorem reference. Separable predual is an intrinsic hypothesis; the historical discussion's chosen Hilbert-space representation is not silently substituted for it.

The candidate quantitative statement allows any unital abelian von Neumann subalgebra \(A\subseteq M\), without requiring a factor or maximality. A conditional expectation onto \(A\) is a positive unital \(A\)-bimodular projection \(M\to A\); normal means ultraweakly continuous. In the report's singular convention, composing the singular expectation \(F\) with normal positive functionals on \(A\) yields singular functionals on \(M\), whose normal component is zero. If both a normal expectation \(E\) and a singular expectation \(F\) exist, the asserted proposition is
\[
 \|E-F\|=2,\qquad
 \sup_{\chi\in\widehat A}\operatorname{diam}\operatorname{Ext}_M(\chi)=2.
\]
The first norm is the operator norm of a map \(M\to A\); fiber diameter is \(\sup_{\varphi,\psi\in\operatorname{Ext}_M(\chi)}\|\varphi-\psi\|_{M^*}\). This is a supremum over pure states and does not assert its attainment by a single specified character. The report supplies no construction of the expectations or full proof.
\subsection{Short English statement}
Determine when every pure state on a MASA has exactly one extension to its factor. The supplied later report gives a separable-predual classification and a norm-two obstruction from normal and singular expectations, while the original 2014 status statements and the arbitrary-predual problem remain separately qualified.
\subsection{Sources}
\begin{itemize}
\item Local UnsolvedMath registry, permanent ID \(20000543\), source code \texttt{catalogue entry}; AIM analysis corpus, \texttt{aim-analysis-notes.json}, notes[74]. Attribution: Ulam AI / UnsolvedMath Contributors, CC BY 4.0.
\item AIM \emph{Beyond Kadison--Singer: paving and consequences}, Question 2; workshop December 1--5, 2014, compilation November 30, 2014: \url{https://docs.google.com/document/d/1XRw0WyG5PshqboCp99tovCRD1GCK0jXZiNLH_nw618c/edit}. The supplied verification reports checking a plain-text export and finding intact wording; this draft claims no new inspection.
\item The source names MSS, the original Kadison--Singer paper, ``Bernard'' in an opening talk, and Popa, Akemann and Sherman for historical discussion. No full citations or expansion of the unidentified Bernard are supplied. Original remarks and literature fields are empty.
\item Registry dated August 14, 2026: statement exact, status \texttt{partially\_solved}, classification \texttt{PARTIAL-PROGRESS}; attempt 1 labelled partial result, rigor proved, novelty confidence low. The later classification and expectation-distance proposition are machine-generated research aids, not peer reviewed, requiring independent verification before citation; their primary references are absent.
\end{itemize}
\end{document}
